\documentclass[10pt]{article}
\usepackage[margin=1in]{geometry}
\usepackage{amsmath,amssymb,amsthm,hyperref}
\newtheorem{theorem}{Theorem}
\newtheorem{lemma}{Lemma}
\title{MathDB 367324: exact boundary values\ for all nonmultiples of three}
\author{Principia Math In Progress}
\date{18 September 2026}
\begin{document}
\maketitle

\section*{Task 1: Complete corrected classification}
Let $M(z)=1+zM(z)+z^2M(z)^2$ and put
\[
 H_N^{(r)}=\det([z^{i+j}]M(z)^r)_{0\le i,j<N},\qquad H_0^{(r)}=1.
\]
Wang and Zhang's Conjecture 7
(\href{https://arxiv.org/abs/2502.21050}{arXiv:2502.21050}, p.~4)
contains sign errors. In particular, its proposed nonzero sum next
to a pair of zero determinants is actually zero. We determine both
summands individually, as well as every other family singled out
by the conjecture, for all exponents and block indices.

\begin{theorem}\label{classification}
Let $r\ge2$, $3\nmid r$, and $n\ge0$. Set
\[
 L=3rn,\qquad \epsilon_r=(-1)^{\binom r2},\qquad
 \eta_r=\begin{cases}1,&r\equiv1\pmod3,\\(-1)^r,&r\equiv2\pmod3.\end{cases}
\]
Then
\begin{align}
 H_L^{(r)}=H_{L+1}^{(r)}&=\epsilon_r^n,\label{c1}\\
 H_{L+r}^{(r)}=H_{L+r+1}^{(r)}&=(-1)^{r-1}\epsilon_r^{n+1},\label{c2}\\
 H_{L+2r}^{(r)}=H_{L+2r+1}^{(r)}&=0,\label{c3}\\
 H_{L+2r-1}^{(r)}=-H_{L+2r+2}^{(r)}
 &=\eta_r\epsilon_r^n[2r(n+1)]^{r-2},\label{c4}\\
 H_{L-1}^{(r)}+H_{L+2}^{(r)}&=-\frac{r(r-3)}2\epsilon_r^n.\label{c5}
\end{align}
For $n=0$, \eqref{c5} uses $H_{-1}^{(r)}=0$; otherwise take $n\ge1$.
For the remaining exponent $r=1$, $H_N^{(1)}=1$ for all $N\ge0$.
\end{theorem}
Thus the source's claimed magnitude $[2r(n+1)]^{r-2}$ is the
magnitude of \emph{each} of the two neighbors in \eqref{c4}, whose
signs are opposite. For example, $r=4,n=0$ gives $H_7^{(4)}=64$
and $H_{10}^{(4)}=-64$. Their sum is zero, not $\pm64$.
Also $H_0^{(4)}=1$ and $H_4^{(4)}=-1$, so the single common sign
printed for all four nonzero boundary values is not valid.
The proof below establishes the complete formulas; the result does
not rely merely on pointing out these discrepancies.

The weighted boundary identities below were established in the
companion solution to
\href{https://github.com/antoshashakov/Principia-Math-In-Progress/blob/mathdb-results/mathdb-open-problems/problems/367323/solution.tex}{MathDB 367323}.
We reproduce their proof and then evaluate the additional minor at
the zero boundaries to obtain the individual values in \eqref{c4}.
For a real parameter $a$, let
$M_a=1+azM_a+z^2M_a^2$, and write its convolution determinants as
$H_N^{(r)}(a)$. Define
\[
 P_j(t)=U_j((t-a)/2),\qquad A(t)=P_{r-1}(t),\qquad
 B_k(t)=U_k\!\left(T_r((t-a)/2)\right).
\]
Here $T_j(\cos\theta)=\cos(j\theta)$ and
$U_j(\cos\theta)=\sin((j+1)\theta)/\sin\theta$; $P_j$ and $B_k$
are monic, of degrees $j$ and $rk$.

\begin{theorem}[Weighted boundary identities]\label{weighted}
For $r\ge2$, $k\ge0$, $N=rk$, and every real or complex $a$,
\begin{align}
 H_N^{(r)}(a)=H_{N+1}^{(r)}(a)&=\epsilon_r^k B_k(0)^{r-1},\label{weighted1}\\
 H_{N-1}^{(r)}(a)+H_{N+2}^{(r)}(a)
 &=\left(r-\binom r2a^2\right)H_N^{(r)}(a).\label{weighted2}
\end{align}
\end{theorem}

\subsection*{The convolution functional and its determinant}

For real $a$, let $\mu_a$ be the translated semicircle measure
\[
 d\mu_a(t)=\frac{\sqrt{4-(t-a)^2}}{2\pi}\,dt\quad(a-2\le t\le a+2).
\]
Its moments are $[z^j]M_a(z)$, and its monic orthogonal polynomials
are $P_j$, with $\int P_iP_j\,d\mu_a=\delta_{ij}$. Both statements
follow by $t=a+2\cos\theta$: the centered moments are the Catalan
numbers and the orthogonality integral is the sine integral.
Set $Q(t)=t^{r-1}A(t)$ and $L(f)=\int f(t)Q(t)\,d\mu_a(t)$.
Then
\begin{equation}\label{functional}
 [z^j]M_a(z)^r=L(t^j).
\end{equation}
Indeed $G(w)=w^{-1}M_a(w^{-1})$ satisfies $G^2=(w-a)G-1$,
so induction gives $G^r=P_{r-1}(w)G-P_{r-2}(w)$. Multiplication
by $w^{r-1}$ and extraction of $w^{-j-1}$ gives \eqref{functional}.
The functional $L$ need not be positive.

We use the following classical Christoffel identity in an algebraic
form. Its evaluation form is Theorem 1 of Krattenthaler,
\href{https://arxiv.org/abs/2101.04225}{arXiv:2101.04225}, with
repeated roots covered by Proposition 5. A proof is included to make
the argument self-contained.

\begin{lemma}\label{remainder}
If $Q$ is monic of even degree $d$, the determinant of the moments
of $Q\,d\mu_a$ of order $m$ is the determinant of the $d\times d$
matrix whose $j$th column is the coefficient vector of
$P_{m+j}\bmod Q$, for $0\le j<d$.
\end{lemma}
\begin{proof}
First suppose $Q$ has distinct roots $\alpha_1,\ldots,\alpha_d$.
Expansion of determinants gives the moment determinant as
\[
 \frac1{m!}\int\cdots\int\Delta(x)^2\prod_{i=1}^m Q(x_i)
 \prod_{i=1}^m d\mu_a(x_i).
\]
The mixed Vandermonde satisfies
$\Delta(x,\alpha)=\Delta(x)\Delta(\alpha)\prod_{i,j}(\alpha_j-x_i)$.
Replace monomials in both Vandermondes by the monic $P_j$.
Orthogonality leaves only columns $0,\ldots,m-1$ for the $x$ rows,
with integral $m!$. Hence the determinant is
\[
 (-1)^{md}\frac{\det(P_{m+j}(\alpha_i))_{1\le i\le d,\,0\le j<d}}
 {\det(\alpha_i^j)_{1\le i\le d,\,0\le j<d}}.
\]
The sign is one. Evaluation of remainders factors the numerator
matrix as the denominator matrix times the coefficient matrix.
The identity extends to repeated roots because both sides are
polynomials in the coefficients of the monic $Q$; polynomial
division introduces no denominators. It also includes $m=0$.
\end{proof}

\subsection*{Exact evaluation at a block boundary}
Put $N=rk$ and $c=N+r-1$. The sine formula gives
\begin{equation}\label{composition}
 P_c=A B_k.
\end{equation}
In the remainder determinant of Lemma~\ref{remainder}, retain its
first $r-1$ columns $P_N,\ldots,P_{N+r-2}$. For each
$1\le j\le r-2$, add column $P_{c-j}$ to column $P_{c+j}$.
The last $r-1$ columns then become
\[
 A B_k S_0,\ldots,A B_k S_{r-2},\qquad
 S_0=1,\quad S_j(t)=2T_j((t-a)/2)\ (j\ge1).
\]
Here $P_{c+j}+P_{c-j}=S_jP_c$, and each $S_j$ is monic of degree $j$.

Use the following monic basis of $\mathbb C[t]/(t^{r-1}A)$:
\[
 P_0,\ldots,P_{r-2},\quad A S_0,\ldots,A S_{r-2}.
\]
Its change-of-basis determinant from the monomials is one. Reduction
modulo $A$ gives an upper block with the first $r-1$ columns and zero
in the last columns. At the simple roots
$\rho_j=a+2\cos(j\pi/r)$ of $A$,
\[
 P_{N+h}(\rho_j)=(-1)^{kj}P_h(\rho_j)\qquad(0\le h\le r-2).
\]
Thus the upper block determinant is
$\prod_{j=1}^{r-1}(-1)^{kj}=\epsilon_r^k$.
The lower right block represents multiplication by $B_k$ in
$\mathbb C[t]/(t^{r-1})$. In its monomial basis this matrix is
triangular, with diagonal $B_k(0)$, so its determinant is
$B_k(0)^{r-1}$. This proves the formula for $H_N$ in
\eqref{weighted1}. No coprimality of $A$ and $t$ was assumed.

\subsection*{The adjacent determinant and the paired sum}
Temporarily choose a real $a$ for which $B_k(0)\ne0$.
This excludes only finitely many values. Equation~\eqref{composition}
gives
\[
 L(t^j B_k)=\int t^{j+r-1}P_{N+r-1}(t)\,d\mu_a(t)=0\qquad(0\le j<N).
\]
Consequently $q_N=B_k$ is the unique monic polynomial of degree $N$
orthogonal to lower degrees for $L$, since $H_N\ne0$.
The elementary identity $U_k(x)^2=\sum_{h=0}^k U_{2h}(x)$ yields
\begin{equation}\label{squared}
 A B_k^2=\sum_{h=0}^k P_{(2h+1)r-1}.
\end{equation}
For $j=0,1,2$, every $h\ge1$ term is orthogonal to $t^{r-1+j}$,
because $r\ge2$. Therefore
\begin{equation}\label{three}
 L(q_N^2)=1,\qquad L(tq_N^2)=ra,\qquad
 L(t^2q_N^2)=r+\binom{r+1}{2}a^2.
\end{equation}
For completeness, the last two values are the coefficients of
$P_{r-1}$ in $t^r$ and $t^{r+1}$. The recurrence
$tP_j=P_{j+1}+aP_j+P_{j-1}$, with $P_{-1}=0$, gives respectively
$ra$ and $\binom{r+1}{2}a^2+r$: in the latter coefficient there
are two level steps or one down step among the otherwise up steps.

The Gram determinant identity
$L(q_N^2)=H_{N+1}/H_N$ now proves $H_{N+1}=H_N$.
In particular $q_{N+1}$ also exists. Suppose first that $N\ge1$.
Let $G_N=(L(t^{i+j}))_{0\le i,j<N}$ and write
\[
 tq_N=q_{N+1}+ra\,q_N+v,\qquad \deg v<N.
\]
For $j<N-1$, $L(t^jv)=0$, whereas $L(t^{N-1}v)=1$.
Thus the coefficient vector of $v$ is $G_N^{-1}e_{N-1}$, and
\[
 L(v^2)=(G_N^{-1})_{N-1,N-1}=\frac{H_{N-1}}{H_N}.
\]
The three summands in the displayed decomposition are mutually
orthogonal. Taking squared norms and using \eqref{three} gives
\[
 r+\binom{r+1}{2}a^2
 =\frac{H_{N+2}}{H_{N+1}}+r^2a^2+\frac{H_{N-1}}{H_N},
\]
which is \eqref{weighted2}. This argument never divides by
$H_{N-1}$ or $H_{N+2}$, which are allowed to vanish.
For $N=0$, the same identity is directly
$H_2=L(t^2)-L(t)^2=r-\binom r2a^2$.
Finally, the asserted identities are polynomial identities in $a$,
so they extend from the generic real values to every complex value,
including $B_k(0)=0$. This proves Theorem~\ref{weighted}.

\subsection*{The missing minor at a zero boundary}
Now put $a=1$, $k=3n+2$, and $N=rk$, with $3\nmid r$.
Write $B=B_k$ and $D=B_{k-1}$. Since $T_r(-1/2)=-1/2$,
\begin{equation}\label{zerovalues}
 B(0)=0,\qquad D(0)=-1,\qquad
 b:=B'(0)=\begin{cases}-2r(n+1),&r\equiv1\pmod3,\\
                         2r(n+1),&r\equiv2\pmod3.\end{cases}
\end{equation}
To verify the derivative, set $t=1+2\cos\theta$. Then
\[
 B(t)=\frac{\sin((k+1)r\theta)}{\sin(r\theta)},\qquad
 \theta\big|_{t=0}=\frac{2\pi}{3},\qquad
 \frac{d\theta}{dt}\bigg|_{t=0}=-\frac1{\sqrt3}.
\]
At this point the numerator is zero and its cosine is one, while
$\sin(2\pi r/3)$ equals $\sqrt3/2$ or $-\sqrt3/2$.
This proves \eqref{zerovalues} and, in particular, $b\ne0$.

Apply Lemma~\ref{remainder} to $H_{N-1}^{(r)}$. Its columns are
\[
 P_{N-1},P_N,\ldots,P_{N+2r-4}.
\]
Move the first column past the next $r-1$ columns, contributing the
sign $(-1)^{r-1}$. As before, use
$P_{c+j}+P_{c-j}=S_jP_c$, now for $1\le j\le r-3$, where
$c=N+r-1$. Since $P_{N-1}=A D$, the reordered, transformed columns are
\[
 P_N,\ldots,P_{N+r-2},\quad A D,\quad
 A B S_0,\ldots,A B S_{r-3}.
\]
For $r=2$ the last list and the corresponding column additions are
empty. In the same monic quotient basis used above, the upper block
again has determinant $\epsilon_r^k$. The lower block is represented
modulo $t^{r-1}$ by
\[
 D,\quad B S_0,\ldots,B S_{r-3}.
\]
Replacing the monic $S_j$ by monomials uses column operations of
determinant one. Because $B(t)=bt+O(t^2)$, the resulting matrix is
triangular with diagonal $D(0),b,\ldots,b$. Hence
\begin{equation}\label{minor}
 H_{N-1}^{(r)}=(-1)^{r-1}\epsilon_r^k D(0)b^{r-2}
              =(-1)^r\epsilon_r^k b^{r-2}.
\end{equation}
Since $H_N=H_{N+1}=0$, the polynomial identity
\eqref{weighted2} gives $H_{N+2}=-H_{N-1}$, with no division by
these zero determinants. For $r\equiv1\pmod3$, the signs from
$(-1)^r$ and $b^{r-2}$ cancel. For $r\equiv2\pmod3$, they give
$(-1)^r$. Also $\epsilon_r^k=\epsilon_r^n$. These observations prove
\eqref{c4} with exactly the stated $\eta_r$.

\subsection*{The other boundary values}
At $a=1$ and $3\nmid r$, the sequence $B_k(0)=U_k(-1/2)$ repeats
$1,-1,0$ with period three. Substitute $k=3n,3n+1,3n+2$ in
\eqref{weighted1} to obtain \eqref{c1}--\eqref{c3}.
Substitution of $k=3n$ into \eqref{weighted2} gives \eqref{c5}.
Finally, when $r=1$, the original semicircle moments have monic
orthogonal polynomials of squared norm one, so every Hankel
determinant is one. This completes the classification and deals
with that separately known exponent explicitly.

\section*{Task 2: Closed formulas inside each block}
Find direct formulas for the residue classes not listed in
Theorem~\ref{classification}, refining the general polynomial-block theorem.

\section*{Task 3: Rank and nullspaces at the zero boundaries}
Determine the exact ranks and explicit nullspace bases of the two
singular Hankel matrices at indices $3rn+2r$ and $3rn+2r+1$.

\section*{Task 4: Least possible block lengths}
Determine the smallest block length for which all residue classes
are polynomial up to a parity factor, and classify exceptional exponents.

\section*{Task 5: Combinatorial cancellation and boundary signs}
Construct an involution on path families that explains both the two
successive zeros and the opposite neighboring signs in \eqref{c4}.
\end{document}
